%%%%%%%%%%  scan idx 609  |  folio 596  %%%%%%%%%%
\section*{1. Énoncé du théorème principal et applications}

\textbf{Théorème 1.1.} \uline{Soient} \(S\) \uline{un schéma noethérien,} \(X\) \uline{et} \(Y\) \uline{deux} \(S\)\uline{-schémas propres,} \(f:X\to Y\) \uline{un morphisme. On suppose} \(S\) intègre \uline{et} \(f\) surjectif. \uline{Alors il existe un ouvert non-vide} \(V\) \uline{de} \(S\) \uline{tel que le morphisme}
\[
  f^*:\underline{\Pic}_{Y/S}|V\longrightarrow \underline{\Pic}_{X/S}|V
\]
\uline{entre les foncteurs de Picard (localisés pour la topologie} f.p.p.f) \uline{soit représentable par un morphisme quasi-affine de présentation finie.}

\textbf{Corollaire 1.2.} \uline{Soit} \(X\) \uline{un schéma propre sur un schéma noethérien} \(S\)\uline{. On suppose} \(S\) intègre. \uline{Alors il existe un ouvert non vide} \(V\) \uline{de} \(S\) \uline{tel que le foncteur de Picard} \(\Pic_{X/S}|V\) \uline{soit représentable par un} \(V\)\uline{-schéma en groupes somme d'ouverts quasi-projectifs de type fini, donc a fortiori séparé sur} \(V\).

En effet rappelons le fait suivant (FGA N\textsuperscript{o}232)~: soient \(S\) un schéma, \(X\) un \(S\)-schéma projectif, de présentation finie et intègre sur \(S\) (c'est-à-dire plat à fibres géométriquement intègres). Alors le foncteur de Picard est représenté par un \(S\)-schéma en groupes \(\underline{\Pic}_{X/S}\), réunion croissante d'ouverts quasi-projectifs de présentation finie.

La considération des «~polynômes de Hilbert~» permet ici de décomposer \(\underline{\Pic}_{X/S}\) en somme d'ouverts \(\underline{\Pic}^q_{X/S}\), \(q\in\mathbb Q[t]\). Dans l'exposé suivant on démontre sous l'hypothèse supplémentaire \(S\) noethérien que les \(\underline{\Pic}^q_{X/S}\) sont de type fini sur \(S\). On sait alors que 1.2 est vrai dans le cas où \(X\) est en plus projectif et intègre sur \(S\); on va en déduire le cas général en utisant le

\textbf{Lemme 1.3.} \uline{Soit} \(X\) \uline{de type fini sur} \(S\) \uline{noethérien intègre. Alors il existe un ouvert non vide} \(V\) \uline{de} \(S\)\uline{, un schéma intègre} \(S'\) \uline{et un morphisme fini et}
